\chapter{Środowisko eksperymentalne i proces ekstrakcji cech}
Część eksperymentalna niniejszej pracy zrealizowana jest w całości w języku \textit{Python}. Środowisko eksperymentalne można natomiast podzielić na trzy części. Pierwszą z nich jest ekstrakcja cech z wybranego modelu bazowego. Cechy te wyodrębniane są z modelu przy użyciu wytrenowanego wcześniej autoenkodera. Kolejnym etapem jest podzielnie ów znalezionych cech na grupy semantyczne, z których każda służyć może jako osobna dziedzina. Ostatnim etapem jest wykorzystanie wydzielonych wcześniej dziedzin do zmodyfikowania modelu bazowego. Pierwszą zbadaną możliwością jest pruning modelu pierwotnego w celu stworzenia modelu dziedzinowego. Drugim eksperymentem jest natomiast budowa modelu typu mieszaniny ekspertów który wybiera odpowiedniego eksperta na podstawie domeny semantycznej. \newline
W niniejszym rozdziale omówiona zostanie pierwsza cześć modułu eksperymentalnego. Opisane zostanie kryterium wyboru modelu bazowego na którym prowadzone były badania, architektura rozwiązania oraz proces treningu modelu autoenkodera. Dodatkowo na każdym etapie wspomniane będą trudności implementacyjne i środowiskowe napotkane w trakcie realizacji rozwiązań, wraz ze sposobami ich unikania i likwidacji. Na końcu krótko omówione zostaną wyniki -- cechy znalezione przez wytrenowany autoenkoder, ich interpretacja oraz analiza.

\section{Wybór modelu bazowego}
Aby rozpocząć eksperymenty związane z ekstrakcją cech, konieczny jest wybór modelu bazowego, na którym realizowane będą wszystkie późniejsze badania. Wybór ten nie jest jednak decyzją jednorazową -- kolejne etapy pracy stawiają przed procesem treningu rzadkiego autoenkodera odmienne wymagania, a dostępne zasoby sprzętowe nie pozwalają zrealizować ich wszystkich na jednym modelu. Z tego powodu w pracy wykorzystane zostały trzy różne modele, z których każdy odpowiada innemu punktowi równowagi pomiędzy pełną kontrolą nad procesem treningu autoenkodera a skalą i jakością modelu bazowego. Model \textit{TinyStories} posłużył jako szybka i niskokosztowa weryfikacja poprawności całego potoku eksperymentalnego; model \textit{Pythia 160M} stanowił właściwą, choć ograniczoną dostępnymi zasobami sprzętowymi, próbę samodzielnego treningu autoenkodera; model \textit{Gemma 3}, wraz z gotowymi autoenkoderami udostępnionymi w ramach \textit{Gemma Scope}, pozwolił natomiast przeprowadzić analizę na modelu klasy produkcyjnej -- niemożliwej do samodzielnego wytrenowania w warunkach niniejszej pracy.
  
  \subsection{Model \textit{TinyStories} jako weryfikacja koncepcji}
  Zanim zaangażowane zostały zasoby potrzebne do treningu autoenkodera na modelu docelowym, konieczne było sprawdzenie, czy sam potok eksperymentalny -- od przechwytywania aktywacji, przez trening rzadkiego autoenkodera, po odczyt wyekstrahowanych cech -- działa poprawnie od początku do końca. Do tego celu wybrany został możliwie najmniejszy i najtańszy obliczeniowo model, na którym błędy implementacyjne dałoby się szybko wykryć i poprawić, zanim zostaną powielone na docelowej, znacznie kosztowniejszej skali. Wyborem tym był model \textit{TinyStories}, przedstawiony przez Ronen Eldan i Yuanzhi Li pracujących w zespole \textit{Microsoft Research} jako dowód na to, że modele rzędu pojedynczych milionów parametrów -- a nawet architektury złożone z zaledwie jednego bloku transformera -- są w stanie generować płynny, spójny i różnorodny tekst, mimo że modele porównywalnej wielkości trenowane na standardowych korpusach (np. \textit{GPT-2} czy \textit{GPT-Neo}) nie wykraczają zwykle poza kilka niespójnych słów \cite{eldan_2023_tinystories_small_language_models}. Konkretny wykorzystany w niniejszej pracy wariant, \textit{TinyStories-1M}, ma jedynie milion parametrów i stanowi najmniejszy z opublikowanych przez autorów punktów kontrolnych -- sam artykuł bada przy tym cały zakres rozmiarów i architektur, nie ograniczając się do tej jednej konfiguracji. \newline
  Trening modelu odbywał się na zbiorze TinyStories -- syntetycznym korpusie krótkich opowiadań wygenerowanych przez modele \textit{GPT-3.5} oraz \textit{GPT-4}, celowo ograniczonych do słownictwa zrozumiałego dla przeciętnego trzy- lub czteroletniego dziecka. To zawężenie słownikowe i tematyczne było warunkiem koniecznym uzyskanego efektu -- pozwoliło autorom oddzielić problem spójności językowej od problemu skali, pokazując, że przy odpowiednio uproszczonej domenie nawet bardzo mały model może opanować podstawy gramatyki i narracji. \newline
  W kontekście niniejszej pracy istotna jest jednak przede wszystkim skrajna wąskość tej domeny. Ograniczone słownictwo i tematyka sprawiają, że autoenkoder wytrenowany na aktywacjach tego modelu operuje na przestrzeni cech nieporównywalnie węższej niż w przypadku docelowych modeli językowych ogólnego przeznaczenia. Z tego powodu wyniki uzyskane na tym etapie mają wyłącznie wartość techniczną -- pozwalają zweryfikować poprawność implementacji całego potoku eksperymentalnego, nie stanowią natomiast podstawy do wnioskowania merytorycznego o sposobie reprezentowania faktów przez model. Po potwierdzeniu poprawności działania potoku na tym uproszczonym modelu możliwe stało się przejście do modelu realistycznej skali, opisanego w kolejnym podrozdziale.

  \subsection{Model \textit{Pythia 160M} jako właściwy trening rzadkiego autoenkodera}
  Mając potwierdzoną poprawność techniczną całego potoku eksperymentalnego, kolejnym krokiem było przejście do modelu na tyle dużego, by trening autoenkodera mógł dać wyniki o rzeczywistej wartości merytorycznej -- pozwalające badać reprezentację faktów, a nie jedynie prostych wzorców językowych. Wyborem tym była rodzina modeli \textit{Pythia} stworzona przez firmę \textit{EleutherAI} oraz opisana w pracy \enquote{\textit{Pythia: A Suite for Analyzing Large Language Models Across Training and Scaling}} \cite{biderman_2023_pythia_suite_analyzing_large}. Główną zaletą tej grupy modeli jest udostępniony przez jej autorów zbiór punktów kontrolnych (ang. \textit{checkpoint}) -- stanów modeli zebranych na różnych etapach treningu. Zestaw ten pozwala na wykorzystanie modelu z dowolnego z $154$ etapów jego nauki -- pozwalając na badanie procesu kształtowania i dystrybucji cech. Ostatecznie element ten nie został jednak jednym z elementów niniejszej pracy. \newline
  Kolejną istotną zaletą modelu \textit{Pythia} jest również jawność użytego do jego nauki zbioru treningowego -- zbioru The Pile \cite{gao_2020_pile_800gb_dataset_diverse}. Ten ważący $800$ GB zestaw danych skonstruowany został z 22 zróżnicowanych, wysokiej jakości podzbiorów. Podzbiory te, opracowane przez zespół \textit{EleutherAI}, obejmują pięć szerokich kategorii tekstu: piśmiennictwo akademickie, literaturę, kod źródłowy (repozytoria GitHub), treści internetowe oraz materiały specjalistyczne, w tym opinie sądowe amerykańskich sądów, zgłoszenia patentowe czy zapisy dyskusji technicznych. Warto również zaznaczyć, że przytłaczająca większość tekstów w tym zbiorze jest w języku angielskim. Z tego też powodu cechy wyodrębnione z modelu \textit{Pythia} będą głównie w języku angielskim lub z nim związane. \newline
  Opisywany model dostępny jest w wariantach od $70$ milionów do $12$ miliardów parametrów. Na różnych etapach pracy testowane były poszczególne rozmiary -- balansując pomiędzy mniejszą złożonością obliczeniową i szybszymi wynikami eksperymentów a większą liczbą cech i lepszą ich jakością przy ekstrakcji, wybór padł na model o rozmiarze $160$ milionów parametrów. Rozmiar ten jest na granicy dostępnych na cele pracy zasobów sprzętowych, a jednocześnie gwarantuje wystarczający rozmiar aby pozwolić na jakościowy trening autoenkodera i ekstrakcję cech. Dlatego więc model \textit{Pythia} został wykorzystany jako główny model eksperymentalny w pracy. \newline
  Mimo wyboru jednego z mniejszych wariantów modelu, czas treningu autoenkodera oraz przestrzeń dyskowa potrzebna na przechowywanie zebranych aktywacji wyraźnie pokazały granice, jakie dostępne zasoby sprzętowe nakładają na samodzielny trening tego typu sieci. Przeprowadzony na tym etapie trening potwierdził wykonalność całego procesu na rzeczywistym modelu językowym, jednocześnie wyznaczając praktyczny sufit skali, powyżej którego dalsze samodzielne eksperymenty przestawały być zasadne -- co bezpośrednio motywuje wykorzystanie w dalszej części pracy gotowego zbioru autoenkoderów.

  \subsection{\textit{Gemma 3} i \textit{Gemma Scope 2}: wykorzystanie gotowych rzadkich autoenkoderów}
  \label{sec:gemma_scope}
  Trening własnego autoenkodera na aktywacjach modelu \textit{Pythia 160M}, opisany w poprzednim podrozdziale, potwierdził, że cały potok eksperymentalny -- od przechwytywania aktywacji, przez trening rzadkiego autoenkodera, po odczyt i interpretację wyekstrahowanych cech -- działa poprawnie i pozwala uzyskać sensowne, interpretowalne wyniki. Powtórzenie tego procesu na modelu rzędu miliardów parametrów wykraczałoby jednak istotnie poza zasoby sprzętowe dostępne na potrzeby niniejszej pracy -- zarówno pod względem czasu potrzebnego na zebranie reprezentatywnego zbioru aktywacji, jak i kosztu samego treningu autoenkodera o odpowiednio dużej przestrzeni cech. Właściwym celem dalszej części pracy jest jednak przetestowanie koncepcji klasteryzacji semantycznej cech oraz budowy modeli domenowych metodami pruningu i architektur typu \textit{Mixture of Experts}. Z tego więc powodu zasadne stało się wykorzystanie gotowego, wcześniej wytrenowanego zestawu autoenkoderów zamiast powtarzania własnego treningu na większą skalę -- pozwala to skoncentrować pozostały budżet czasowy i obliczeniowy pracy na właściwym przedmiocie dalszych rozdziałów, korzystając jednocześnie z cech o jakości i skali nieosiągalnej w warunkach samodzielnego treningu. \newline
  W tym celu wykorzystany został zbiór \textit{Gemma Scope 2} \cite{gemma_scope_2_2025} -- opublikowany w grudniu 2025 roku przez zespół \textit{Google DeepMind}, otwarty zestaw rzadkich autoenkoderów oraz transkoderów (ang. \textit{transcoders}), wytrenowanych na aktywacjach całej rodziny modeli \textit{Gemma 3}, obejmującej warianty od $270$ milionów do $27$ miliardów parametrów. \newline
  Dla każdej warstwy modelu \textit{Gemma 3} udostępniono kilka wariantów autoenkodera, różniących się szerokością słownika cech (m.in. warianty $16$k, $65$k oraz $262$k kierunków) oraz docelowym poziomem rzadkości aktywacji (parametr $L_0$). % TODO: uzupełnić konkretny wybór -- rozmiar modelu Gemma 3, numer/nazwa warstwy, szerokość i L0 wykorzystanego autoenkodera oraz krótkie uzasadnienie tego wyboru
  \newline
  Model \textit{Gemma 3} wraz z odpowiadającym mu zbiorem autoenkoderów z \textit{Gemma Scope 2} stanowi od tego miejsca docelowy obiekt badań niniejszej pracy -- to na nim przeprowadzona zostanie właściwa analiza sąsiedztwa semantycznego cech, próba pruningu i budowy architektury \textit{Mixture of Experts} oraz końcowa walidacja i wizualizacja wyników. Modele \textit{TinyStories} oraz \textit{Pythia}, opisane w poprzednich podrozdziałach, pełnią względem niego funkcję pomocniczą -- odpowiednio jako narzędzie walidacji potoku oraz jako dowód wykonalności samodzielnego treningu autoenkodera na rzeczywistym modelu językowym, zanim ta sama umiejętność została wykorzystana do świadomej rezygnacji z powtarzania treningu na większą skalę.
  %TODO: w przyszłosci warto to zmodyfikować lub sprostować

\section{Konfiguracja potoku analitycznego}
\label{sec:konfiguracja_potoku}
Trening rzadkiego autoenkodera wymaga zbioru wektorów aktywacji wybranej warstwy modelu bazowego. Przygotowanie tych wektorów odbywa się w kilku następujących po sobie etapach: surowy korpus tekstowy jest najpierw dzielony na zdania, tokenizowany i pakowany w sekwencje o stałej długości, następnie przetwarzany przez model, a aktywacje wybranej warstwy są zbierane i przekazywane do procesu treningu autoenkodera. Kolejne podrozdziały szczegółowo opisują poszczególne etapy, a na rysunku \ref{fig:potok_sae} przedstawiono diagram architektury całego potoku.
\begin{figure}[ht]
    \centering
    \begin{tikzpicture}[
        node distance=0.6cm and 0.7cm,
        box/.style={draw, rounded corners=3pt, minimum height=1.0cm,
                    text width=3.4cm, align=center, font=\footnotesize},
        data/.style={box, fill=blue!8},
        proc/.style={box, fill=orange!12},
        arr/.style={->, >=stealth, thick},
    ]
    % --- Rząd 1 ---
    \node[data] (corpus) {Korpus tekstowy};
    \node[proc, right=of corpus] (seg) {Segmentacja zdań};
    \node[proc, right=of seg] (tok) {Tokenizacja};
    
    % --- Rząd 2 ---
    \node[proc, below=0.8cm of tok] (pack) {Pakowanie\\w sekwencje};
    \node[proc, left=of pack] (pad) {Wyrównanie długości\\i maska uwagi};
    \node[data, left=of pad] (npy) {Sekwencje tokenów\\gotowe do przetworzenia};
    
    % --- Rząd 3 ---
    \node[proc, below=0.8cm of npy] (hook) {Ekstrakcja aktywacji\\wybranej warstwy};
    \node[proc, right=of hook] (sae) {Trening\\rzadkiego autoenkodera};
    \node[data, right=of sae] (features) {Wyekstrahowane\\cechy modelu};
    
    % --- Strzałki ---
    \draw[arr] (corpus) -- (seg);
    \draw[arr] (seg) -- (tok);
    \draw[arr] (tok) -- (pack);
    \draw[arr] (pack) -- (pad);
    \draw[arr] (pad) -- (npy);
    \draw[arr] (npy) -- (hook);
    \draw[arr] (hook) -- (sae);
    \draw[arr] (sae) -- (features);
    \end{tikzpicture}
    \caption{Diagram potoku przygotowania danych i treningu rzadkiego autoenkodera. Proces obejmuje wyłącznie ścieżkę ekstrakcji cech; etapy separacji domen semantycznych i budowy modeli ekspertowych opisane są w dalszych rozdziałach.}
    \label{fig:potok_sae}
\end{figure}

  \subsection{Przygotowanie sekwencji danych treningowych}
  \label{sec:przygotowanie_sekwencji}
  W celu przygotowania sekwencji tekstowych użyte zostały dwa opisane wcześniej zbiory danych -- zbiór \textit{TinyStories} dla mniejszego modelu oraz zbiór \textit{The Pile} dla modeli \textit{Pythia} i \textit{Gemma}. Model \textit{TinyStories} wymagał użycia dedykowanego dla niego zbioru tekstów, na których był uczony. Zbiór \textit{The Pile} cechuje się natomiast wystarczającą różnorodnością, aby zastosować go dla obu większych modeli. \newline
  Pierwszym etapem przygotowywania sekwencji jest podział korpusu na zdania za pomocą biblioteki \texttt{Syntok} \cite{leitner_syntok_2022}. Biblioteka ta realizuje lingwistyczną segmentację zdań -- rozpoznaje granice zdań na podstawie znaków interpunkcyjnych i kontekstu, dzięki czemu żadne zdanie nie jest rozcinane w połowie. Podział na poziomie zdań, a nie na poziomie ustalonej liczby znaków czy słów, zapobiega sytuacji, w której model otrzymuje niepełne zdanie, co mogłoby prowadzić do artefaktów w zbieranych aktywacjach. \newline
  Kolejnym krokiem jest tokenizacja wyodrębnionych zdań przy użyciu tokenizatora odpowiadającego danemu modelowi. Każde zdanie zamieniane jest na ciąg tokenów -- liczb całkowitych odpowiadających pozycjom słów w słowniku modelu. Zdania dłuższe niż zadana długość sekwencji są na tym etapie dzielone na fragmenty nieprzekraczające tego limitu. \newline
  Na tym etapie warto dodać, że do wczytywania odpowiadających danym modelom tokenizatorów jak i do ładowania wag i konfiguracji wszystkich modeli bazowych wykorzystywanych w tej i kolejnych częściach środowiska eksperymentalnego użyta została biblioteka \textit{Hugging Face Transformers} \cite{wolf_2020_huggingfaces_transformers_state_of_the_art_natural}. \newline
  Stokenizowane zdania trafiają następnie do etapu pakowania. Algorytm utrzymuje bufor, do którego sekwencyjnie dołączane są kolejne zdania, dopóki ich łączna długość nie przekroczy docelowej długości sekwencji. Gdy dodanie kolejnego zdania spowodowałoby przekroczenie limitu, bieżąca zawartość bufora jest zapisywana jako kompletna sekwencja, a nowe zdanie rozpoczyna kolejny bufor. Fragmenty krótsze niż ustalony próg minimalnej długości (domyślnie $10$ tokenów) są odrzucane, aby uniknąć degeneracji treningu na zbyt krótkich ciągach. \newline
  Po zakończeniu pakowania następuje wyrównanie długości sekwencji. Każda sekwencja uzupełniana jest tokenami paddingu do dokładnie tej samej, docelowej liczby pozycji. Jednocześnie generowana jest binarna maska uwagi (ang. \textit{attention mask}) -- macierz o tym samym kształcie co sekwencje, w której wartość $1$ oznacza prawdziwy token, a $0$ -- token paddingu. Maska ta jest niezbędna, ponieważ w słowniku niektórych modeli identyfikator $0$ odpowiada prawidłowemu tokenowi, a zatem sama wartość numeryczna nie pozwala odróżnić go od paddingu. \newline
  Oprócz właściwych sekwencji procedura generuje plik metadanych zawierający informacje niezbędne do reprodukcji eksperymentu: ścieżki plików źródłowych, wykorzystaną frakcję zbioru, ziarno generatora liczb pseudolosowych (ang.\ \textit{seed}), łączną liczbę tokenów oraz liczbę wygenerowanych sekwencji.

  \subsection{Zbieranie aktywacji stanów ukrytych}
  \label{sec:zbieranie_aktywacji}
  Po przygotowaniu sekwencji tokenów kolejnym etapem potoku jest przechwytywanie wewnętrznych reprezentacji wygenerowanych przez badany model językowy. W procesie przetwarzania przez blok transformera (omówionym w podrozdziale~\ref{sec:architektura_transformer}), każda pozycja w sekwencji tokenów przekształcana jest w wektor stanu ukrytego w strumieniu rezydualnym (ang. \textit{residual stream}). Wybór punktu ekstrakcji ma istotne znaczenie dla analizy -- pośrednie warstwy modelu odpowiadają za reprezentowanie złożonych cech semantycznych oraz wiedzy faktograficznej, zanim zostaną one zagregowane i zmapowane na bezpośrednie prawdopodobieństwa słownika wyjściowego. \newline
  Do przechwytywania wektorów aktywacji wykorzystany został mechanizm punktów przechwytujących (ang. \textit{forward hooks}), udostępniany przez bibliotekę PyTorch \cite{paszke_2019_pytorch_imperative_style_high_performance} -- główny framework głębokiego uczenia stosowany w całym potoku eksperymentalnym. Jest to funkcja rejestrowana na wybranym module sieci, która podczas wykonywania przejścia w przód (ang. \textit{forward pass}) automatycznie przechwytuje i przekazuje wartości tensorów wyjściowych bez konieczności ingerencji w strukturę modelu. Rozwiązanie to pozwala na bieżący dostęp do stanów ukrytych bez naruszania oryginalnej architektury. \newline
  Przechwycone aktywacje obejmują stany ukryte dla wszystkich pozycji w sekwencji. Przed przekazaniem ich do dalszych etapów potoku stosuje się filtrowanie pozycji odpowiadających tokenom paddingu przy użyciu binarnej maski uwagi. W efekcie otrzymuje się wyłącznie wektory aktywacji odpowiadające rzeczywistym tokenom tekstu, o wymiarze równym szerokości stanu ukrytego modelu ($d_{\mathrm{model}}$), które stanowią bezpośrednie wejściowe dane treningowe dla rzadkiego autoenkodera. \newline
  Zarówno sekwencje tokenów, jak i przechwycone aktywacje przechowywane są w postaci tablic biblioteki NumPy \cite{harris_2020_array_programming_with_numpy}, z wykorzystaniem mechanizmu odwzorowania plików w pamięci (ang. \textit{memory-mapped files}). Mechanizm ten polega na mechanizmie sprawiającym, że system operacyjny traktuje plik na dysku jak rozszerzenie pamięci operacyjnej -- poszczególne fragmenty tablicy ładowane są do RAM dopiero w momencie rzeczywistego dostępu, a nieużywane strony mogą być zwolnione i ponownie odczytane z dysku w razie potrzeby. Dzięki temu program może indeksować tablicę o rozmiarze przekraczającym dostępną pamięć operacyjną tak, jakby była ona w całości załadowana, co w kontekście niniejszej pracy było koniecznym zabiegiem ze względu na rozmiar zbiorów aktywacji sięgający kilkudziesięciu lub nawet kilkuset gigabajtów.
  
  \subsection{Przetwarzanie strumieniowe -- ograniczenia sprzętowe i modyfikacja kodu}
  \label{sec:przetwarzanie_strumieniowe}
  W pierwotnym podejściu proces ekstrakcji aktywacji podzielony był na dwa osobne etapy: wygenerowanie i zapisanie pełnego zbioru aktywacji na dysku, a następnie ładowanie tak przygotowanego pliku do treningu autoenkodera. Wraz ze wzrostem skali eksperymentów i przejściem do większych modeli językowych podejście to ujawniło krytyczne ograniczenie sprzętowe. Przechowywanie wektorów aktywacji dla milionów tokenów wymagało dziesiątek gigabajtów przestrzeni dyskowej oraz znacznego narzutu pamięci operacyjnej na operacje wejścia-wyjścia, stanowiąc bariery trudne do pokonania w warunkach ograniczonego budżetu sprzętowego. \newline
  Aby rozwiązać ten problem, potok eksperymentalny został zmodyfikowany i oparty na przetwarzaniu strumieniowym. Zamiast tworzenia pojedynczego pliku z aktywacjami na dysku, kod wykorzystuje generator przetwarzający dane w niewielkich porcjach. W każdej iteracji model bazowy wykonuje przejście w przód tylko dla wybranego fragmentu sekwencji, przechwycone aktywacje są natychmiast przekazywane do pętli treningowej autoenkodera w pamięci operacyjnej, a po wykonaniu aktualizacji wag autoenkodera pamięć jest zwalniana. Pozwala to na trening autoenkodera bez konieczności trwałego składowania aktywacji na dysku. \newline  
  Istotnym elementem architektury strumieniowej jest również mechanizm wznowienia pracy połączony z punktami kontrolnymi (ang. \textit{checkpoints}). Zapisywany stan obejmuje nie tylko wagi autoenkodera i stan optymalizatora, ale także kursor określający wskaźnik ostatnio przetworzonej sekwencji oraz stany generatorów liczb pseudolosowych. W przypadku przerwania procesu treningowego wykonanie może zostać wznowione od dokładnie tego samego miejsca bez utraty spójności danych i bez potrzeby ponownego przeliczania wcześniejszych fragmentów korpusu.

\section{Trening Top-K SAE}
\label{sec:trening_topk_sae}
Proces treningu SAE dotyczy przede wszystkim modelu \textit{Pythia 160M} -- w przypadku \textit{Gemma 3} wykorzystywany jest gotowy zestaw autoenkoderów \textit{Gemma Scope 2}, opisany w kolejnym podrozdziale. Zastosowana implementacja realizuje architekturę Top-K SAE \cite{gao_2024_scaling_evaluating_sparse_autoencoders}, opisaną szczegółowo w podrozdziale \ref{sec:sae_warianty}, i działa bezpośrednio na strumieniu aktywacji opisanym w poprzednich podrozdziałach. \newline
Autoenkoder przyjmuje na wejściu wektor aktywacji $x \in \mathbb{R}^{d_{\mathrm{model}}}$. Przed właściwym kodowaniem od wektora tego odejmowany jest wyuczony wektor przesunięcia dekodera $b_{\mathrm{dec}}$, centrujący dane wokół zera:
\[
\tilde{x} = x - b_{\mathrm{dec}}.
\]
Tak wycentrowany wektor przekształcany jest liniowym enkoderem, po czym poddawany nieliniowości ReLU:
\[
z = \mathrm{ReLU}(W_{\mathrm{enc}} \tilde{x} + b_{\mathrm{enc}}).
\]
Z otrzymanego wektora pobudzeń $z \in \mathbb{R}^{d_{\mathrm{SAE}}}$, o wymiarze $d_{\mathrm{SAE}} = d_{\mathrm{model}} \cdot r$, gdzie $r$ oznacza współczynnik rozszerzenia (ang. \textit{expansion factor}), zachowywanych jest jedynie $k$ największych wartości, pozostałe zaś zerowane -- zgodnie z mechanizmem selekcji cech opisanym w podrozdziale \ref{sec:sae_warianty}:
\[
\hat{z} = \mathrm{TopK}_k(z).
\]
Rekonstrukcja aktywacji wejściowej powstaje jako liniowa kombinacja wyselekcjonowanych kierunków słownika, do której z powrotem dodawany jest ten sam wektor przesunięcia:
\[
\hat{x} = W_{\mathrm{dec}} \hat{z} + b_{\mathrm{dec}}.
\]
Istotne jest przy tym, że macierz dekodera $W_{\mathrm{dec}}$ nie jest transpozycją macierzy enkodera $W_{\mathrm{enc}}$, lecz osobnym, niezależnie inicjalizowanym i trenowanym parametrem. Ten sam wektor $b_{\mathrm{dec}}$ pełni natomiast podwójną rolę, służąc zarówno do centrowania danych przed enkodowaniem, jak i do przesunięcia rekonstrukcji z powrotem do oryginalnej przestrzeni aktywacji. \newline
Funkcją straty jest wyłącznie błąd średniokwadratowy (ang. \textit{Mean Squared Error}, MSE) pomiędzy rekonstrukcją a oryginalną aktywacją:
\[
\mathcal{L} = \lVert \hat{x} - x \rVert_2^2.
\]
W odróżnieniu od klasycznego SAE trenowanego karą $L1$, funkcja straty nie zawiera żadnego dodatkowego członu wymuszającego rzadkość -- rzadkość reprezentacji jest tu konsekwencją strukturalną mechanizmu Top-K. Parametry sieci optymalizowane są algorytmem AdamW bez regularyzacji wagowej (ang. \textit{weight decay}), z obcinaniem normy gradientu do wartości maksymalnej równej $1$ w celu ustabilizowania treningu. Tempo uczenia jest dostosowywane przez mechanizm wyżarzania kosinusowego (ang. \textit{cosine annealing}), malejącego w sposób ciągły od wartości początkowej do zadanej wartości minimalnej wraz z kolejnymi krokami optymalizacji. \newline
Podczas treningu dla każdej cechy słownika prowadzony jest licznik liczby jej wystąpień wśród $k$ aktywnych kierunków w danym kroku optymalizacji. Na tej podstawie po zakończeniu treningu wyliczany jest odsetek cech martwych (ang. \textit{dead features}) -- kierunków, które w trakcie całego treningu ani razu nie znalazły się wśród $k$ najsilniej aktywnych. \newline
Warstwą, z której pobierane są aktywacje modelu \textit{Pythia 160M}, jest warstwa o numerze $6$ -- warstwa środkowa spośród $12$ warstw tego modelu. Wymiar wejściowy autoenkodera $d_{\mathrm{model}}$ nie jest ustalany ręcznie, lecz odczytywany automatycznie z konfiguracji wczytanego modelu bazowego (dla \textit{Pythia 160M} wynosi on $768$), a jawnie podana wartość jest w kodzie porównywana z tą odczytaną, co zapobiega niezauważonemu niedopasowaniu wymiarów pomiędzy autoenkoderem a modelem, na którego aktywacjach jest on trenowany. \newline

  \subsection{Adaptacja gotowych SAE (\textit{Gemma Scope}) do potoku}
  \label{sec:adaptacja_gemma_scope}
  W przypadku modelu \textit{Gemma 3} etap treningu autoenkodera jest zastąpiony załadowaniem gotowych wag z zestawu \textit{Gemma Scope 2} \cite{gemma_scope_2_2025}, opisanego w podrozdziale~\ref{sec:gemma_scope}. Wagi udostępnione są w repozytorium \textit{Hugging Face} i ładowane za pośrednictwem biblioteki \textit{SAELens} \cite{bloom_2024_sae_training_codebase} -- otwartego narzędzia do treningu i analizy rzadkich autoenkoderów, które udostępnia API do wczytywania autoenkoderów z różnych źródeł. \newline
  Istotną różnicą architektoniczną jest to, że autoenkodery z \textit{Gemma Scope 2} wykorzystują architekturę JumpReLU \cite{rajamanoharan_2024b_jumprelu_sparse_autoencoders}, a nie Top-K stosowane we własnym treningu autoenkodera na bazie modelu \textit{Pythia}. Autoenkoder JumpReLU stosuje uczony, indywidualny próg aktywacji dla każdej cechy zamiast globalnego parametru $k$. W konsekwencji liczba aktywnych cech nie jest stała, lecz zależy od konkretnej aktywacji wejściowej. Z tego powodu moduł analizy cech dla modelu \textit{Gemma 3} został zrealizowany jako osobny skrypt, korzystający z enkodera udostępnianego przez \textit{SAELens}, przy zachowaniu identycznej logiki zbierania statystyk i generowania raportu. \newline
  Do analizy wybrano wariant autoenkodera oznaczony jako \texttt{layer\_9\_width\_16k\_l0\_medium}, odpowiadający warstwie $9$ -- środkowej spośród $18$ warstw modelu \textit{Gemma 3 270M}. Słownik tego autoenkodera zawiera $16\,384$ cechy, co stanowi kompromis pomiędzy rozdzielczością analizy a kosztem obliczeniowym. Punkt ekstrakcji aktywacji jest natomiast dla \textit{Gemma Scope 2} analogiczny do punktu wykorzystywanego w potoku dla modelu \textit{Pythia}. \newline
  Pozostała część potoku analitycznego -- przygotowanie sekwencji tokenów z korpusu \textit{The Pile}, przechwytywanie stanów ukrytych przez punkt przechwytujący zarejestrowany na wybranej warstwie modelu oraz strumieniowe przetwarzanie aktywacji -- została przeniesiona bezpośrednio z istniejącej infrastruktury opisanej w podrozdziale~\ref{sec:konfiguracja_potoku}. Analiza cech -- obejmująca statystyki aktywacji, przykłady kontekstowe oraz przybliżoną analizę \textit{logit lens} -- przebiega identycznie jak dla modelu \textit{Pythia} i generuje raporty w tym samym formacie tekstowym i strukturalnym.

\section{Interpretacja i wstępna analiza cech}
\label{sec:interpretacja_cech}
Podstawową informacją charakteryzującą każdą cechę są jej statystyki aktywacji zebrane na całym korpusie treningowym. Dla każdej cechy $i$ wyznaczane są: maksymalna wartość aktywacji $\hat{z}_i^{\max}$, średnia wartość aktywacji po warunkowaniu na tokeny, dla których $\hat{z}_i > 0$ (tzw. \textit{mean activation on active tokens}), oraz częstość aktywacji (ang. \textit{activation frequency}) -- odsetek tokenów ze zbioru treningowego, dla których cecha jest aktywna, tzn. $\hat{z}_i > 0$. Miara ta jest wskaźnikiem rzadkości cechy: cechy aktywujące się na ułamku procenta tokenów reprezentują koncepcje wąskie i wyspecjalizowane, podczas gdy cechy o wysokiej częstości aktywacji mogą odpowiadać ogólnym właściwościom składniowym lub statystycznym tekstu. \newline
Na podstawie zebranych statystyk możliwe jest wykrycie cech martwych (ang. \textit{dead features}) -- wektorów, które przez cały trening nie zostały ani razu aktywowane powyżej progu. Cechy te nie niosą żadnej informacji semantycznej i są wykluczone z dalszej analizy. Analogicznie identyfikowane są cechy nadaktywne (ang. \textit{ultra-high density features}) -- kierunki aktywujące się na blisko każdym tokenie, które zazwyczaj kodują niskopoziomowe statystyki tekstu niezwiązane ze znaczeniem. \newline
Zasadniczą metodą interpretacji cech jest przegląd tekstowych przykładów ich aktywacji. Dla każdej analizowanej cechy z korpusu treningowego wyodrębniane są fragmenty tekstu, w których dana cecha uzyskała najwyższe wartości aktywacji. Każdy taki fragment obejmuje kilkanaście tokenów poprzedzających token maksymalnej aktywacji oraz kilka tokenów po nim, co pozwala na odczytanie kontekstu. \newline
Przeglądając zbiór takich przykładów -- zazwyczaj od kilku do kilkudziesięciu najwyżej aktywowanych fragmentów -- możliwe jest zidentyfikowanie wspólnego mianownika: konkretnego słowa, klasy, kontekstu tematycznego lub wzorca składniowego, który jest wspólny dla wszystkich wystąpień danej cechy. Właśnie ta spójność tematyczna stanowi definicję interpretowalności cechy stosowaną w niniejszej pracy: cechę uznaje się za interpretowalną, jeśli analizujący jest w stanie sformułować zwięzłe, konsekwentne z przykładami wyjaśnienie tego, co cecha wykrywa. \newline
Opisana procedura wymaga oceny i opisu przez człowieka. Zapewnia ona jednak informację, że zidentyfikowany kierunek w przestrzeni aktywacji odpowiada jakiejś spójnej, interpretowalnej koncepcji. Wyniki tej analizy przeprowadzonej dla wybranych cech modelu \textit{Pythia} zestawione zostały w~aneksie~\ref{appendix:katalog_cech_pythia}, który stanowi katalog wybranych cech które zostały wybrane i zinterpretowane przez autora wraz ze wskazaniem tokenów je aktywujących oraz przykładami wystąpień tych cech w tekście. \newline
Wyniki tej analizy należy interpretować ostrożnie: ponieważ SAE działa na pośredniej warstwie modelu, a nie na bezpośrednim wejściu warstwy odwzorowania na słownik, uzyskiwane tokeny mają jedynie charakter orientacyjny. Niemniej zestawienie listy tokenów z \textit{logit lens} z obserwowanymi przykładami wystąpień pozwala z dobrą skutecznością potwierdzić znaczenie danej cechy. \newline
Tak skonstruowane opisy i statystyki każdej z cech są następnie używane jako reprezentacja semantyczna w procesie klasteryzacji opisanym w rozdziale~\ref{chap:klasteryzacja}, gdzie cechy o zbliżonych właściwościach są grupowane w domeny semantyczne służące jako podstawa do budowy modeli ekspertowych.